% =====================================================================
\section*{Техника безопасности и рабочее место}
\addcontentsline{toc}{section}{Техника безопасности и рабочее место}
\markboth{Техника безопасности и рабочее место}{}

Опыты этого курса идут на безопасном напряжении \SIrange{3.3}{5}{V}. Ударить током оно не может, но
испортить плату, обжечь палец о перегретую деталь или повредить аккумулятор можно за секунду.
Правила ниже короткие. Прочитайте их до первой сборки и возвращайтесь к ним перед каждым новым уроком.

\subsection*{Десять правил}

\begin{enumerate}
  \item \textbf{Сначала отключите питание, потом меняйте схему.} Выньте USB-кабель, пересоберите,
        проверьте по рисунку и только после этого включайте.
  \item \textbf{На выводы GPIO --- не больше \SI{3.3}{V}.} Контакт 5V --- это питание, а не сигнал.
  \item \textbf{Не замыкайте 3V3 и 5V на GND.} Оголённые концы проводов и металлические предметы на столе ---
        частая причина короткого замыкания.
  \item \textbf{Деталь греется или пахнет --- сразу отключите питание.} Потом ищите ошибку: чаще
        всего перепутаны плюс и минус.
  \item \textbf{Светодиод подключайте через резистор, а мотор и реле --- через ключ или драйвер}
        (уроки 1 и 16).
  \item \textbf{Сетевое напряжение \SI{220}{V} в учебных опытах не используется.} Никогда не подключайте
        к макетной плате провода от розетки и не разбирайте блоки питания.
  \item \textbf{Литиевые аккумуляторы требуют особой осторожности:} не замыкать, не прокалывать, не
        заряжать без платы защиты и без присмотра (урок 27). Вздувшийся аккумулятор отдайте педагогу.
  \item \textbf{Паяльник --- только с разрешения педагога.} Жало нагрето выше \SI{300}{\celsius}: кладите
        паяльник только на подставку, работайте в проветриваемом помещении, после пайки мойте руки.
  \item \textbf{Берегите глаза.} Откусывая выводы деталей, придерживайте их и направляйте вниз, от себя
        и от соседей. Не смотрите на яркие светодиоды в упор.
  \item \textbf{Порядок на столе.} Никакой еды и напитков рядом с электроникой; в конце занятия отключите
        питание и разложите детали по ячейкам.
\end{enumerate}

\begin{caution}
Статическое электричество способно повредить микросхему незаметно для вас. Перед работой коснитесь
заземлённого металлического предмета, не трогайте контакты модуля камеры и чипа пальцами и держите
платы за края.
\end{caution}

\subsection*{Если что-то пошло не так}

\begin{itemize}
  \item \textbf{Дым, искры, запах гари} --- отключите питание, отойдите, позовите педагога. Не пытайтесь
        потушить или разобрать устройство самостоятельно.
  \item \textbf{Ожог} --- подержите обожжённое место под прохладной водой и сообщите педагогу.
  \item \textbf{Плата перестала определяться компьютером} --- отключите всё от её выводов и проверьте
        кабель; чаще всего виновато короткое замыкание в схеме, а не «сгоревшая» плата.
\end{itemize}

% =====================================================================
\section*{Макетная плата, схема и мультиметр}
\addcontentsline{toc}{section}{Макетная плата, схема и мультиметр}
\markboth{Макетная плата, схема и мультиметр}{}

Три инструмента понадобятся в каждом уроке: макетная плата, на которой собирают схему без пайки,
принципиальная схема, по которой её собирают, и мультиметр, которым проверяют результат.

\subsection*{Как устроена макетная плата}

Под отверстиями макетной платы спрятаны металлические полоски. Они соединяют отверстия группами, и это
нужно запомнить один раз.

\begin{center}
\begin{tikzpicture}[x=4mm,y=4mm,font=\sffamily\scriptsize]
  \fill[gray!8,draw=gray!60,rounded corners=2pt] (-1.2,-1.4) rectangle (30.2,15.4);
  % шины питания: соединены вдоль
  \foreach \y/\c in {14.4/esp,13.4/accent,0.6/esp,-0.4/accent}{
    \draw[\c,line width=2pt,opacity=0.35] (0,\y) -- (29,\y);
    \foreach \x in {0,...,29} \fill[black!70] (\x,\y) circle (1pt);
  }
  % рабочее поле: соединены пятёрками поперёк
  \foreach \x in {0,...,29}{
    \draw[okgreen,line width=2pt,opacity=0.3] (\x,7.6) -- (\x,11.6);
    \draw[okgreen,line width=2pt,opacity=0.3] (\x,2.4) -- (\x,6.4);
    \foreach \y in {2.4,3.4,4.4,5.4,6.4,7.6,8.6,9.6,10.6,11.6} \fill[black!70] (\x,\y) circle (1pt);
  }
  \node[anchor=west,text=esp] at (30.8,14.4) {шина «+»: вдоль всей платы};
  \node[anchor=west,text=accent] at (30.8,13.4) {шина «$-$» (GND)};
  \node[anchor=west,text=okgreen!70!black,align=left] at (30.8,9.6) {пять отверстий в столбике\\ соединены между собой};
  \node[anchor=west,align=left] at (30.8,7.0) {канавка: связи нет};
  \draw[thin] (30.6,7.0) -- (29.6,7.0);
  \node[anchor=west,text=okgreen!70!black,align=left] at (30.8,4.4) {нижние пятёрки ---\\ отдельные от верхних};
\end{tikzpicture}
\end{center}

\begin{itemize}
  \item \textbf{Шины питания} по краям соединены \emph{вдоль} платы. К красной подключают 3V3, к синей --- GND.
  \item \textbf{Рабочее поле} соединено \emph{поперёк}, столбиками по пять отверстий. Если вставить две
        ножки детали в один столбик, они окажутся замкнуты.
  \item \textbf{Канавка} посередине разделяет верхние и нижние пятёрки. Плату ESP32-S3 и микросхемы ставят
        над канавкой, чтобы ножки левого и правого ряда не соединялись.
\end{itemize}

\begin{caution}
На многих длинных макетных платах шины питания разорваны посередине: левая и правая половины не
соединены. Если половина схемы «не получает питание», поставьте перемычки между половинами шин или
проверьте их мультиметром в режиме прозвонки.
\end{caution}

\subsection*{Как читать схему}

Принципиальная схема показывает не внешний вид деталей, а их соединения. Линия --- это провод, точка на
пересечении линий --- соединение. Чаще всего в книге встречаются такие обозначения:

\begin{center}
\begin{circuitikz}[font=\sffamily\scriptsize]
  \draw (0,0) to[R] (2,0);            \node at (1,-0.8) {резистор};
  \draw (2.7,0) to[led] (4.7,0);      \node at (3.7,-0.8) {светодиод};
  \draw (5.4,0) to[push button] (7.4,0); \node at (6.4,-0.8) {кнопка};
  \draw (8.1,0) to[C] (10.1,0);       \node at (9.1,-0.8) {конденсатор};
  \draw (10.8,0) to[pR] (12.8,0);     \node at (11.8,-0.8) {потенциометр};
  \draw (13.8,-0.2) -- (13.8,0.3) node[vcc]{}; \node at (13.8,-0.8) {питание};
  \draw (15.2,0.5) -- (15.2,0.1) node[ground]{}; \node at (15.2,-0.8) {земля};
\end{circuitikz}
\end{center}

Все значки «земля» на схеме соединены между собой, даже если линия между ними не нарисована. То же
относится к одинаково подписанным питаниям (3,3 В). У светодиода треугольник показывает направление
тока: от анода (длинная ножка) к катоду.

Номинал резистора можно прочитать по цветным полоскам: первые две --- цифры, третья --- число нулей.

\begin{center}\small
\begin{tabular}{@{}lllll@{}}
\toprule
чёрный --- 0 & коричневый --- 1 & красный --- 2 & оранжевый --- 3 & жёлтый --- 4 \\
зелёный --- 5 & синий --- 6 & фиолетовый --- 7 & серый --- 8 & белый --- 9 \\
\bottomrule
\end{tabular}
\end{center}

Например, «красный, красный, коричневый» --- это 2, 2 и один ноль: \SI{220}{\ohm}. «Коричневый, чёрный,
оранжевый» --- 1, 0 и три нуля: \SI{10}{k\ohm}. Золотая или серебряная полоска в конце показывает
точность. Если сомневаетесь, измерьте резистор мультиметром.

\subsection*{Мультиметр: четыре измерения}

\begin{center}\small
\begin{tabularx}{\textwidth}{@{}l l X@{}}
\toprule
Что измеряем & Режим & Как подключать щупы \\
\midrule
Напряжение & $V$ с прямой чертой (DC), предел 20 В & \emph{параллельно}: чёрный щуп на GND, красный --- в точку схемы; питание включено \\
Целостность цепи & прозвонка (значок звука или диода) & к двум концам провода или дорожки; \textbf{питание выключено} \\
Сопротивление & $\Omega$ & к выводам детали, вынутой из схемы; \textbf{питание выключено} \\
Ток & $mA$ или $A$ (DC) & \emph{последовательно}, в разрыв цепи; красный щуп иногда переставляют в другое гнездо \\
\bottomrule
\end{tabularx}
\end{center}

\begin{caution}
В режиме измерения тока мультиметр ведёт себя как кусок провода. Если в этом режиме коснуться щупами
плюса и минуса источника, получится короткое замыкание. После измерения тока сразу возвращайте
переключатель и красный щуп в положение для напряжения.
\end{caution}

Первое, что делают при любой неполадке, --- измеряют питание: между 3V3 и GND должно быть около
\SI{3.3}{V}, между 5V и GND --- около \SI{5}{V}. Если напряжения нет, искать ошибку в программе бесполезно.

\begin{tip}
Нет платы или деталей под рукой? Многие опыты первых уроков можно собрать в браузерном симуляторе
Wokwi (\url{https://wokwi.com}): там есть ESP32-S3, светодиоды, кнопки, дисплеи и Serial Monitor.
Симулятор не заменит настоящую сборку, но помогает проверить программу и схему до занятия.
\end{tip}

\begin{tasks}
\begin{enumerate}
\item \lv{1} Сформулируйте вслух три правила безопасности, которые нарушать нельзя, и объясните соседу, почему.
\item \lv{1} Найдите на макетной плате ряды питания и покажите, какие отверстия в одном столбце соединены между собой.
\item \lv{2} Определите по полоскам номиналы трёх резисторов и проверьте их мультиметром. Совпали ли значения? Почему допустимо небольшое расхождение?
\item \lv{2} Измерьте напряжения 3V3--GND и 5V--GND при питании от USB и запишите результаты: они пригодятся при поиске неисправностей.
\item \lv{3} Соберите на макетной плате цепь «3V3 --- резистор --- светодиод --- GND», нарисуйте её схему и объясните, в какую сторону включён светодиод.
\item \lv{4} Мультиметром в режиме прозвонки найдите, не разорваны ли шины питания посередине вашей платы, и составьте карту: какие участки соединены, а какие нет.
\end{enumerate}
\end{tasks}
